\documentclass[conference]{IEEEtran}
\IEEEoverridecommandlockouts

\usepackage{amsmath}
\usepackage{amssymb}
\usepackage{booktabs}
\usepackage{graphicx}
\usepackage{url}

\title{Query-to-Grasp: Language-Queryable RGB-D Retrieval, Multi-View Memory, and H200-Scale Simulated Manipulation in ManiSkill}

\author{\IEEEauthorblockN{Anonymous Authors}
\IEEEauthorblockA{Paper draft scaffold for internal development}}

\begin{document}
\maketitle

\begin{abstract}
Language-conditioned manipulation systems need more than a 2D language
grounding result: they need a 3D target that is geometrically consistent,
diagnosable under ambiguity, and usable by an action executor. We present
Query-to-Grasp, a modular H200-validated ManiSkill prototype for
language-queryable RGB-D target retrieval, multi-view object memory,
confidence-aware re-observation diagnostics, and opt-in simulated manipulation
evaluation. The system connects
GroundingDINO detections, optional CLIP reranking, RGB-D lifting, object-centric
3D fusion, deterministic target selection, and a scripted top-down ManiSkill
pick and pick-place controllers. The default executor remains a safe
placeholder, while simulated executors report downstream \texttt{pick\_success},
\texttt{place\_success}, and task-level \texttt{task\_success} separately. On \texttt{PickCube-v1}, fused memory grasp
points yield full-query multi-view and closed-loop simulated pick success of
\texttt{1.0000}. On \texttt{StackCube-v1}, task-aware guarded multi-view
targets reach \texttt{0.6200} tabletop and \texttt{0.5200} closed-loop
pick-only success over 50 seeds, while oracle object poses reach
\texttt{0.9400} pick success and \texttt{0.8800} privileged pick-place task
success. The accepted query-pick/oracle-place bridge reaches
\texttt{0.7200} single-view task success, \texttt{0.5200} tabletop task
success, and \texttt{0.4800} closed-loop task success with an oracle cubeB
destination. These results
support Query-to-Grasp as a reproducible diagnostic baseline for studying the
gap between open-vocabulary RGB-D retrieval, executable simulated targets, and
task-level manipulation outcomes.
\end{abstract}

\section{Introduction}

Open-vocabulary detectors can localize language-described objects in images,
but robotic manipulation requires a stronger interface. A manipulation system
must lift detections into a consistent 3D frame, merge evidence across views,
decide whether the selected target is sufficiently supported, and expose when
an apparently confident target is not yet graspable. Query-to-Grasp studies
this intermediate layer between language grounding and action execution.

The system is deliberately modular. It parses a natural-language query, obtains
2D proposals with GroundingDINO~\cite{liu2023groundingdino}, optionally reranks
candidates with CLIP~\cite{radford2021learning}, lifts RGB-D observations into
world coordinates, fuses multi-view detections into object memory, and selects
a query-matching 3D target. A placeholder pick executor remains the default so
retrieval and uncertainty metrics can be benchmarked without overstating
control capability. An opt-in simulated top-down executor then tests whether
the selected target can drive ManiSkill actions.

This paper makes three claims. First, careful RGB-D geometry matters: after an
OpenCV-to-OpenGL camera-convention correction, multi-view memory becomes
compact enough for deterministic fusion and traceable target selection. Second,
downstream pick success depends on preserving a grasp target separately from a
semantic object center: fused \texttt{memory\_grasp\_world\_xyz} targets lift
\texttt{PickCube-v1} multi-view and closed-loop pick success to \texttt{1.0000}.
Third, broader tasks expose useful limitations: \texttt{StackCube-v1} supports
pick-only compatibility, oracle pick-place establishes a privileged placement
upper bound, and the accepted bridge shows that query-derived cubeA targets can
support placement when cubeB is oracle-provided.

No real-robot execution is claimed. The controller is scripted, not learned,
and StackCube results do not claim stack placement completion. The contribution
is instead a benchmarkable systems scaffold with per-run artifacts, target
source diagnostics, and reproducible limitations for language-conditioned
RGB-D manipulation research.

\section{Related Work}

\subsection{Open-Vocabulary Grounding}

Open-vocabulary vision-language models provide reusable perception primitives
for language-conditioned robotics. CLIP~\cite{radford2021learning} aligns
image and text embeddings, while GroundingDINO~\cite{liu2023groundingdino}
grounds language prompts into 2D detections. Query-to-Grasp uses these models
as off-the-shelf proposal and reranking modules. Our focus is not improving
grounding accuracy; it is measuring what happens after a language-conditioned
2D proposal must become a 3D target and then a simulated grasp command.

\subsection{Language-Conditioned Manipulation}

Language-conditioned manipulation systems connect instructions to actions
through policies, planners, or affordance models. CLIPort~\cite{shridhar2022cliport}
combines language-conditioned semantic features with spatial action maps,
PerAct~\cite{shridhar2022peract} uses voxelized transformer policies,
SayCan~\cite{ahn2022saycan} grounds language-model plans in robot affordances,
and VIMA~\cite{jiang2023vima} studies multimodal prompt-conditioned
manipulation. RT-1~\cite{brohan2022rt1} and RT-2~\cite{brohan2023rt2} scale
robot policies with large datasets and vision-language representations.
Query-to-Grasp is narrower: it does not learn a manipulation policy, but it
provides a transparent retrieval-to-target layer that such systems often need.

\subsection{RGB-D Memory and Simulated Manipulation Benchmarks}

RGB-D manipulation systems often rely on explicit geometry, object-centric
memory, or multi-view fusion to convert perception into action state. ManiSkill
benchmarks~\cite{gu2023maniskill2} provide controlled simulated environments
for evaluating manipulation pipelines. Query-to-Grasp uses ManiSkill both as
an RGB-D observation source and as the execution substrate for a scripted
top-down pick controller. The key design choice is to keep semantic memory
separate from grasp-target memory so retrieval confidence and downstream
control can be diagnosed independently.

\subsection{Active and Re-Observation Diagnostics}

Active perception and next-best-view methods use additional observations to
resolve uncertainty before acting. Query-to-Grasp implements a diagnostic,
rule-based re-observation loop rather than a learned view planner. The policy
requests an extra virtual view when confidence gaps, view support, point
support, or geometry signals are weak. This exposes a useful systems lesson:
closed-loop re-observation can reduce uncertainty metrics without improving
pick success when the extra view is absorbed by the wrong object or shifts the
effective target source.

\section{Method}

\subsection{Single-View Query-to-3D Pipeline}

Given a language query, a deterministic parser extracts a target phrase,
attributes, and conservative relation fields. The perception layer detects
query-conditioned 2D boxes, optionally applies CLIP reranking, and lifts
selected RGB-D pixels into 3D. The lifted semantic center is stored as
\texttt{world\_xyz}. For simulated pick experiments, the system can also compute
an additive grasp target, \texttt{grasp\_world\_xyz}, from local RGB-D geometry.
Semantic retrieval remains unchanged when this grasp point is available.

\subsection{Multi-View Object Memory}

The multi-view path captures a virtual \texttt{tabletop\_3} preset and adds
per-view detections to object memory. Each memory object stores label votes,
view support, confidence terms, point counts, semantic centers, and optional
grasp observations. Semantic centers support retrieval and target selection.
Grasp observations support downstream simulated pick execution. This separation
allows Table~\ref{tab:target-source} to compare \texttt{selected\_object\_world\_xyz},
\texttt{memory\_grasp\_world\_xyz}, task-guarded semantic targets, and
\texttt{oracle\_object\_pose} without changing the retrieval score.

\subsection{Re-Observation Policy}

The re-observation policy is rule-based and diagnostic. It triggers on low
overall confidence, ambiguous top candidates, insufficient view support,
insufficient 3D point support, or geometry outliers. A closed-loop benchmark
can capture one extra view, merge it into memory, and rerun target selection.
Continuity guards keep selected-object association stable when appropriate,
while absorber-aware checks prevent forcing continuity when a competing object
also absorbed the extra view.

\subsection{Simulated Pick and Pick-Place Execution}

The default executor returns a structured placeholder result. The opt-in
\texttt{sim\_topdown} executor uses ManiSkill \texttt{pd\_ee\_delta\_pos}
actions to move the TCP above the target, descend, close the gripper, and lift.
The opt-in \texttt{sim\_pick\_place} executor extends the sequence with a
scripted move-to-place, release, and settle stage. Both executors report
\texttt{grasp\_attempted}, \texttt{pick\_success}, \texttt{place\_success},
\texttt{task\_success}, stage diagnostics, final TCP position, and simulator
info. \texttt{pick\_success} measures task-specific grasp or lift evidence,
while \texttt{task\_success} is the raw ManiSkill task success flag.

\begin{figure}[t]
\centering
\fbox{\parbox{0.92\linewidth}{Fig.~1 placeholder. Pipeline diagram from
\texttt{outputs/paper\_figure\_pack\_latest/notes/implemented\_architecture.md}.}}
\caption{Query-to-Grasp pipeline: language query, RGB-D grounding, 3D lifting,
multi-view memory, re-observation diagnostics, and optional simulated pick
execution.}
\label{fig:pipeline}
\end{figure}

\section{Experimental Setup}

Experiments run in ManiSkill with RGB-D observations. H200 HF/ManiSkill runs
are the source of paper-scale claims; local and Colab paths are retained for
smoke testing. The primary benchmark is \texttt{PickCube-v1}. \texttt{StackCube-v1}
is the cross-task bridge benchmark: accepted rows test query-driven cubeA
picking, oracle rows test privileged cubeA-to-cubeB placement, and the accepted
query-pick/oracle-place bridge tests query-derived cubeA targets with an oracle
cubeB placement target.

Metrics are split into retrieval, uncertainty, and control groups. Retrieval
metrics include detection count, 3D target availability, memory size, selected
confidence, and target source. Re-observation metrics include trigger rate,
closed-loop resolution, and still-needed rate. Control metrics include
\texttt{grasp\_attempted\_rate}, \texttt{pick\_success\_rate},
\texttt{place\_success\_rate}, \texttt{task\_success\_rate}, and stage counts.

The main paper artifacts are collected in \texttt{outputs/paper\_figure\_pack\_latest}.
Table~\ref{tab:target-source} summarizes the final target-source results.
Additional reports include the single-view ablation table, corrected geometry
reports, full PickCube grasp comparison, StackCube expanded failure report,
oracle target-source comparison, and the accepted query-driven StackCube
placement bridge.

\begin{table*}[t]
\centering
\caption{Target-source and simulated manipulation results used by the current
paper draft. The StackCube bridge uses query-derived cubeA targets and a
privileged oracle cubeB placement target; oracle pick-place is fully
privileged.}
\label{tab:target-source}
\begin{tabular}{llrrrrl}
\toprule
Benchmark & Environment & Runs & Failed & Pick success & Task success & Effective target source \\
\midrule
Full multi-view/closed-loop & \texttt{PickCube-v1} & 55 & 0 & 1.0000 & 0.1455 & \texttt{memory\_grasp\_world\_xyz} \\
Expanded tabletop & \texttt{StackCube-v1} & 50 & 0 & 0.6200 & 0.0000 & \texttt{task\_guard\_selected\_object\_world\_xyz} \\
Expanded closed-loop & \texttt{StackCube-v1} & 50 & 0 & 0.5200 & 0.0000 & \texttt{task\_guard\_selected\_object\_world\_xyz} \\
Query-pick/oracle-place single-view & \texttt{StackCube-v1} & 50 & 0 & 0.8800 & 0.7200 & query target\,$\rightarrow$\,\texttt{oracle\_cubeB\_pose} \\
Query-pick/oracle-place tabletop & \texttt{StackCube-v1} & 50 & 0 & 0.6200 & 0.5200 & query target\,$\rightarrow$\,\texttt{oracle\_cubeB\_pose} \\
Query-pick/oracle-place closed-loop & \texttt{StackCube-v1} & 50 & 0 & 0.5200 & 0.4800 & query target\,$\rightarrow$\,\texttt{oracle\_cubeB\_pose} \\
Oracle pick & \texttt{PickCube-v1} & 50 & 0 & 1.0000 & 0.0400 & \texttt{oracle\_object\_pose} \\
Oracle pick & \texttt{StackCube-v1} & 50 & 0 & 0.9400 & 0.0000 & \texttt{oracle\_object\_pose} \\
Oracle pick-place & \texttt{StackCube-v1} & 50 & 0 & 0.9400 & 0.8800 & \texttt{oracle\_cubeA\_pose}\,$\rightarrow$\,\texttt{oracle\_cubeB\_pose} \\
\bottomrule
\end{tabular}
\end{table*}

\section{Results}

\subsection{Grounding and Reranking Baseline}

The single-view baseline establishes that the language-to-3D chain is stable
before adding multi-view memory. Exact PickCube queries produce a 3D target in
all validated runs. Ambiguity queries increase candidate counts but do not
show CLIP as the main source of improvement in the current benchmarks: top-1
changes are rare because candidate multiplicity is low.

\subsection{Camera Geometry Enables Memory}

The OpenCV-to-OpenGL camera-convention correction is the first critical
systems result. After the correction, same-label cross-view spread drops from
\texttt{1.0693 m} to \texttt{0.0518 m}, and memory fragmentation drops from
\texttt{3.3333} to \texttt{1.3333} objects per run. Fig.~\ref{fig:geometry}
should use the geometry reports from the paper pack to visualize this change.

\begin{figure}[t]
\centering
\fbox{\parbox{0.92\linewidth}{Fig.~2 placeholder. Camera-convention and
cross-view geometry reports from
\texttt{outputs/paper\_figure\_pack\_latest/geometry/}.}}
\caption{Correcting the camera convention turns multi-view coordinates from a
dominant failure mode into a usable object-memory signal.}
\label{fig:geometry}
\end{figure}

\subsection{Fused Grasp Points Enable PickCube Picking}

Semantic fused centers are useful for retrieval, but they are not always good
pick targets. Propagating per-view grasp points into fused memory and using
\texttt{memory\_grasp\_world\_xyz} for refined multi-view simulated picking
raises compact and full-query PickCube pick success to \texttt{1.0000}. This is
the strongest positive result in the current paper: the retrieval target, fused
memory, and scripted controller can agree well enough to execute successful
simulated picks across the full ambiguity set.

\subsection{StackCube Exposes Cross-Task Limits}

StackCube validates that the pipeline can transfer to a second ManiSkill task
for pick-only control, but it also exposes task-specific target-source
preferences. A StackCube guard uses the selected object's semantic fused center
as the effective refined multi-view target, yielding \texttt{0.6200} tabletop
and \texttt{0.5200} closed-loop pick success over 50 seeds. The oracle baseline
reaches \texttt{0.9400}, which indicates that many remaining StackCube failures
are target-source and association limitations rather than an absolute inability
of the scripted controller to pick cubeA.

\subsection{Closed-Loop Re-Observation Is Diagnostic, Not Always Beneficial}

Closed-loop re-observation reduces uncertainty diagnostics on PickCube and in
the accepted absorber-aware policy path, but it does not improve StackCube pick
success. The expanded StackCube failure report attributes failures to wrong
fused grasp observations, memory fragmentation or low support, controller or
contact residuals, and closed-loop third-object absorption. This supports a
conservative interpretation: re-observation is useful for diagnosis and can
resolve selection uncertainty, but it is not automatically a grasp-success
improvement.

\section{Limitations}

Query-to-Grasp is a simulated systems prototype. It does not claim real-robot
deployment, a learned controller, or full long-horizon manipulation. The
\texttt{sim\_topdown} executor is a scripted top-down pick policy, and
\texttt{sim\_pick\_place} is a scripted placement probe rather than a learned
controller. Accepted query-driven StackCube rows remain pick-only until the
query-pick/oracle-place bridge is enabled with a privileged destination target.
Oracle pick-place rows are privileged upper bounds, not deployable
language-conditioned stacking.

The detector backend is used off the shelf. The current results do not show
CLIP reranking as a primary improvement source. The query parser handles simple
attributes and conservative relation fields, not unrestricted language
reasoning. Multi-view re-observation is implemented with virtual views in
ManiSkill, not physical camera motion.

\section{Conclusion}

Query-to-Grasp connects language-queryable RGB-D perception to multi-view 3D
memory, uncertainty-aware re-observation, and opt-in simulated pick and
pick-place execution.
The current system shows that careful camera geometry and separate grasp-point
memory are enough to turn PickCube from a retrieval-only benchmark into a
strong simulated pick benchmark. StackCube demonstrates useful cross-task
transfer while exposing target-source, placement-target, and closed-loop
association limits. These results provide a reproducible H200-scale basis for
an IROS/ICRA-style simulated systems paper and a clear roadmap toward broader
task success.

\appendices
\section{Artifact and Command Notes}

The current paper pack is generated with:
\begin{verbatim}
python scripts/build_paper_figure_pack.py \
  --output-dir outputs/paper_figure_pack_latest \
  --skip-missing
\end{verbatim}

The LaTeX structure check is:
\begin{verbatim}
python scripts/check_paper_latex.py \
  --tex paper/main.tex \
  --bib paper/references.bib
\end{verbatim}

The main benchmark artifacts referenced by this draft are:
\begin{itemize}
  \item \texttt{outputs/h200\_60071\_multiview\_memory\_grasp\_point\_full\_ambiguity\_seed01234}
  \item \texttt{outputs/h200\_60071\_stackcube\_task\_guard\_expanded\_seed0\_49}
  \item \texttt{outputs/h200\_60071\_oracle\_pick\_ablation}
  \item \texttt{outputs/h200\_60071\_oracle\_stackcube\_place\_seed0\_49}
  \item \texttt{outputs/h200\_60071\_query\_stackcube\_place\_bridge\_seed0\_49}
  \item \texttt{outputs/paper\_figure\_pack\_latest}
\end{itemize}

\bibliographystyle{IEEEtran}
\bibliography{references}

\end{document}
